\chapter*{ЗАКЛЮЧЕНИЕ}
\addcontentsline{toc}{chapter}{ЗАКЛЮЧЕНИЕ}

В результате выполнения данной работы цель достигнута: разработан метод поддержки принятия решений при выполнении пошаговой инструкции, заданной в графовом формате. Решены все задачи:

\begin{enumerate}
    \item проведен анализ предметной области систем поддержки принятия решений, сопровождающих выполнение набора предписаний человеком, выполнен сравнительный анализ методов принятия решений, описан существующий графовый формат описания инструкций, на примере выбранного домена формализованы классы возможных состояний, событий и связанных с ними действий;
    \item разработан метод поддержки принятия решений в рамках выполнения пошаговой инструкции, заданной в графовом формате, описаны основные особенности предлагаемого метода, сформулированы ограничения предметной области, изложены ключевые этапы метода в виде функциональной модели и схем алгоритмов, выделены структуры данных, необходимые для реализации метода;
    \item обоснован выбор средств программной реализации, разработано программное обеспечение, реализующее представленный метод, выполнено его тестирование, описано взаимодействие пользователя с программным обеспечением;
    \item описана исследовательская процедура, составлен набор сценариев развития событий для каждой решаемой задачи в рамках выбранного домена, включая различные виды нештатного развития ситуации, выполнена параметризация метода согласно описанной исследовательской процедуре, даны рекомендации о применимости метода; также проведена оценка качества результатов предложенного метода.
\end{enumerate}

В результате проведенного исследования начальные предположения о том, что предложенная модификация метода анализа иерархий с группировкой альтернатив позволяет существенно сократить временные затраты эксперта по сравнению с классическим подходом, а экспертная настройка весов критериев в режиме отладки повышает качество выдаваемых рекомендаций, подтвердились.

% \newpage

% В ходе данной научно-исследовательской работы был проведен анализ предметной области систем поддержки принятия решений, сопровождающих выполнение набора предписаний человеком, выполнен сравнительный анализ методов принятия решений, описан графовый формат описания инструкций, формализованы классы состояний, событий и действий на примере выбранного домена, выполнена формализация постановки задачи в виде функциональной модели. Разработанный метод сочетает в себе гибридное графовое представление инструкции, аппарат метода анализа иерархий для ранжирования рекомендаций, механизм обработки аварийных сценариев с поддержкой рекурсии и возможность экспертной настройки весов критериев в режиме реального времени.

% Таким образом, все задачи для достижения цели данной работы были решены, и цель работы — разработка метода поддержки принятия решений при выполнении пошаговой инструкции, заданной в графовом формате — была достигнута. Разработанный метод и программное обеспечение могут быть рекомендованы для использования в кулинарных приложениях, обучающих системах, промышленных инструкциях и других областях, связанных с пошаговым выполнением предписаний.